\begin{frame}{일정과 마일스톤 M1~M4}
  \small
  \begin{tabular}{@{}p{0.6cm}p{7.2cm}p{5.8cm}@{}}
    \toprule
    주 & 해야 할 일 & 마일스톤(제출) \\
    \midrule
    1 & 데이터 후보 2~3개 조사, \textbf{정형/비정형 판정}, 팀 회의,
      \textbf{문제 정의 문서}(행·라벨·지표·베이스라인·``왜 DL'' 초안) &
      \textbf{M1} 문제 정의 + 데이터 선택 + ``왜 DL'' 1페이지 \\
    2 & 3분할(Ch06.3) 스캐폴드, \textbf{train으로만 fit한 전처리},
      \textbf{1개 신경망 엔드투엔드 확인}, \textbf{GBDT 베이스라인} &
      \textbf{M2} 중간코드 리뷰 (코드 + 베이스라인 숫자 각 1줄) \\
    3 & 나머지 기법 적용, val/CV(Ch06.2)로 튜닝, \textbf{test 최종
      측정(한 번만)}, 학습 곡선 & \textbf{M3} 학습 곡선 1장 + 최종
      test 숫자 + GBDT 비교 \\
    4 & 수식적 정당화, 보고서·슬라이드, \textbf{발표(5~7분, 마지막에
      ``왜 DL'')}, \textbf{peer-review 2팀} & \textbf{M4} 보고서 +
      발표 + 리뷰 2건 \\
    \bottomrule
  \end{tabular}
  \vskip0.5em
  \begin{itemize}
    \item 4주(각 주 = 수업 3회). M1·M2는 개인별 선정 $\to$ 회의에서
      팀 전체가 확정
    \item M2가 \textbf{베이스라인 숫자 한 줄}을 요구: 봐야 ``신경망을
      만들었다''고 말할 수 있는지 판정 가능
    \item M3가 \textbf{학습 곡선 한 장}을 요구: 최종 숫자만으론 에포크
      를 어떻게 정했는지 검증 불가 — val 곡선이 실제로 쓰였는지
      그림에서 드러남. Block B의 x축은 \textbf{에포크}(8.1의
      $n\_estimators$와 다름)
  \end{itemize}
\end{frame}
